\documentclass[12pt]{article}
\usepackage[margin=1in]{geometry}
\usepackage[T1]{fontenc}
\usepackage{xcolor}
\definecolor{garnet}{HTML}{73000A}

\setlength{\parindent}{0pt}
\setlength{\parskip}{0.3em}

\begin{document}
\noindent
Department of Mechanical Engineering\\
University of South Carolina\\
Columbia, SC 29208\\
\today
\vspace{0.5em}
\noindent

{\color{garnet}\rule{\textwidth}{0.5pt}}

\vspace{0.3em}

\noindent
Dear Members of the Evaluation Committee:

I recommend Samuel Cancilla for the SC Space Grant Undergraduate Research Award. As his faculty advisor in the Integrated Multiphysics and Systems Engineering Laboratory (iMSEL) since March 2025, I have directly supervised his research on adaptive radar tracking for autonomous maritime perception.

Samuel came to iMSEL with prior research experience: he spent the summer of 2023 in Dr.\ Austin Downey's Adaptive Real-Time Systems Laboratory (ARTS-Lab), where he designed pulse generator and power amplifier circuits for piezoelectric vibration damping on instrumented test beams---learning to scope hardware experiments, interpret frequency-domain measurements, and iterate on circuit prototypes under advisor feedback. He is now co-author on a paper submitted to AIAA titled ``Surface-Object Detection and Tracking from Commercial X-Band Marine Radar Using Adaptive ST-DBSCAN,'' which presents an algorithm he developed and validated from the ground up.

Samuel planned and led multi-site radar data collection campaigns across three South Carolina lakes---Murray, Monticello, and Greenwood---accumulating over 30\,hours of X-band radar recordings under varying sea states. He then traveled to Charleston Harbor for 8\,additional hours of multi-modal sensor data collection in an operational maritime environment. During the Greenwood campaign, unexpected afternoon wind gusts shifted lake clutter conditions mid-collection; rather than discarding the session, Samuel revised the collection grid on-site to capture both calm and rough-water regimes, turning a logistics problem into a richer dataset. These campaigns required coordinating boat logistics, managing equipment in field conditions, and adapting collection protocols when conditions changed. In my eight years supervising undergraduate researchers, only two others have independently organized multi-site field campaigns of this scope.

Samuel faced a core technical challenge: standard clustering algorithms require hand-tuned parameters that break down when radar clutter changes between environments. He chose to build an adaptive algorithm that self-tunes using local point-cloud statistics, then implemented the full real-time tracking pipeline in Rust to meet the throughput constraint of our 48\,RPM Furuno NXT marine radar on an embedded NVIDIA Jetson Orin platform. When early Rust profiling revealed that memory allocation in the clustering loop caused frame drops, he restructured the pipeline around pre-allocated buffers---a systems-level decision he made independently before bringing it to our group meeting. He then designed and ran a systematic ablation study comparing his adaptive tracker against fixed-parameter baselines, demonstrating measurable improvements in track continuity.

His proposed research extends this work in two directions: generalizing the adaptive parameter estimation to include temporal self-tuning, and training a classifier to distinguish small boats, large ships, and clutter from cluster-level features. I have reviewed his methodology. The 85\% classification accuracy target is grounded in published results: Vivone et~al.\ (2021) achieved 82\% on a comparable maritime radar feature set with a simpler classifier, and Samuel's feature space adds track-level statistics that were unavailable in that work. He completed the full adaptive ST-DBSCAN pipeline---from algorithm design through embedded deployment and ablation testing---in 10\,weeks, and the proposed extensions are narrower in scope, making the 14-week timeline realistic.

Samuel himself identified the NASA relevance of his work during a lab presentation, drawing the connection between maritime radar perception in GPS-denied harbor environments and analogous challenges facing autonomous systems on lunar or planetary surfaces. The underlying problem---adaptive spatiotemporal clustering of noisy point clouds without human-tuned parameters---applies directly to SAR and LiDAR sensors used in NASA Earth observation and surface exploration missions.

Samuel also serves as an Army National Guard 09R SMP Cadet through ROTC and as Assistant Editor of The Daily Gamecock's opinion section. He maintains a 3.5\,GPA while balancing military obligations, editorial deadlines, and research fieldwork. He prepares equipment checklists before each field campaign, maintains detailed logs during collection, and debriefs afterward to identify protocol improvements---the kind of operational discipline that keeps multi-day data collection on track.

I give Samuel Cancilla my strongest recommendation for this award.

\vspace{0.8em}

\noindent
Sincerely,

\vspace{1.5em}

\noindent
Dr.\ Yi Wang\\
Associate Professor\\
Department of Mechanical Engineering\\
University of South Carolina\\
\textit{yiwang@cec.sc.edu}

\end{document}
